\section{2024-09-13 - Interferometer}

\startsectionlevel[title={Video -
Gravitationswellen},reference={video---gravitationswellen}]

\startitemize[packed]
\item
  Gravitationswellen stauchen und strecken materie minimal
\item
  Solche wellen werden unter anderem durch die Kollision von schwarzen
  Löchern verursacht
\stopitemize

\stopsectionlevel

\startsectionlevel[title={Das
Michelson-Interferometer},reference={das-michelson-interferometer}]

\startframedtask
{\bf Aufgabe:} Erklären Sie mithilfe der S. 173 unter Erstellung einer
Skizze das Funktionsprinzip eines Michelson Interferometers und die
Messung kleiner Längenänderungen damit.

\stopframedtask

\margintext{
\startplacefigure[title={Michelson-Interferometer},reference={fig:ex:inferferometer},location={here}]
\externalfigure[assets/physik/wellenoptik/Interferometer.svg][marginwidth]
\stopplacefigure
}

Durch die Trennung und die wieder Zusammenführung des Laser (\in{Abbildung}[fig:ex:inferferometer]) kann mit der
Interferenz eine beeinflussende Kraft mit sehr hoher präzision
festgestelt werden.



\startframedexample
{\bf Gangunterschied}

Der Gangunterschied $\delta s$ ist der Unteschied der Laufwege der
beiden Wellen bei Interferenz

\stopframedexample

\startframedexample
Aufgabe bei Leifiph

\stopframedexample

\stopsectionlevel


\stopsectionlevel
